{\small\textbf{DNB's framework (Diego, 1:15) [10:10–11:25]}}\par\smallskip
That compass is DNB's 2022 framework, with four phases.
\begin{itemize}\setlength\itemsep{0pt}
\item After a crisis, the buffer is released.
\item In \textbf{normal} times, it should be at \textbf{two percent}, built up one point a year.
\item Only when risks are clearly elevated does it go above two percent.
\item When a crisis hits, it is released again.
\end{itemize}
\par\smallskip
\textcolor{navy}{\textbf{[def]}} This is a "positive neutral" buffer: a buffer that is already filled when risks are neither high nor low.
\par\smallskip
How did DNB get to two percent? From history.
\begin{itemize}\setlength\itemsep{0pt}
\item In past crises, Dutch banks' peak accumulated losses were about \textbf{12 billion euros}.
\item A two-percent buffer is about \textbf{six billion} of releasable capital, and releasing it could support up to \textbf{150 billion} of lending.
\item \textcolor{navy}{\textbf{[def]}} Decisions are taken by "guided discretion": a dashboard of indicators informs, but does not dictate.
\end{itemize}
\par\smallskip
The Netherlands is not alone: seventeen jurisdictions now run such a framework, and Sweden, the UK and Poland also target two percent.
